\section{Appendix}

In this section we gather some facts most of which can already be found in 
\cite{SchikorraIntBoundFrac11} in slightly different versions. The main aim is to prove 
Lemma~\ref{la:tangentialpart} and \ref{lem:CriticalTerm} which both rely on 
quasi-locality of the Riesz potential $\lapms{s}$.  Afterwards, we give an easy proof of the iteration lemma needed
to deduce Dirichlet growth. 

We will mainly deal with functions belonging to the Schwartz space $\mathcal{S}(\R)$
of rapidly decreasing smooth functions $\R\to\R$. The statements carry over to more general situations by suitable approximation arguments.

\subsection{Quasi-locality}

The essential tool apart from Sobolev inequalities is the following 
quantitative version of the \emph{quasi-locality} of the fractional Laplacian 
and the Riesz potential.

\begin{lemma}[Quasi-locality]\label{lem:QuasiLocality}
 Let $p_1,p_2,q_1,q_2 \in [1, \infty]$, where we assume that $q_2 =1$ if $p_2=1$, $s \in (-1,1)$
 and $\Omega_1, \Omega_2$ be disjoint domains with
 $d:= \dist(\Omega_1, \Omega_2) >0$ and with positive and finite Lebesgue measure. 
 Then, for any $f\in\mathcal S(\R)$, 
 \begin{equation*}
  \|\lap^{\frac{s}{2}} (f \chi_{\Omega_2})\|_{(p_1,q_1),\Omega_1} 
  \aleq  d^{-1-s} |\Omega_1|^{1/p_1} |\Omega_2|^{1-1/p_2} \|f\|_{(p_2,q_2), \Omega_2}
 \end{equation*}
  where we set
 \[
  \lap^{\frac{s}{2}} := \begin{cases}
                         \laps{s} \quad &\mbox{if $s > 0$,}\\
			  \id \quad &\mbox{if $s = 0$,}\\
			  \lapms{\abs{s}} \quad &\mbox{if $s < 0$.}\\
                        \end{cases}
 \]
\end{lemma}

Note, that $\|\cdot\|_{(1,q_1)}$, $\|\cdot\|_{(\infty,q_\infty)}$ are only considered
in the inequalities that follow for $q_1=1$ and $q_\infty=\infty$.

\begin{proof}
 For $\tilde k_s (z) = \frac 1 {|z|^{1+s}} \chi_{\R\setminus B_d} (z)$
 and $\supp f\subset\Omega_{2}$ we have
 for all $x \in \Omega_1$
 \begin{equation*}
  \lap^{\frac{s}{2}} f(x) = c_{s}(\tilde k_s \ast f ) (x)
 \end{equation*}
 and hence
 \begin{equation*}
 |\lap^{\frac{s}{2}} f (x)| \aleq \|f\|_{1,\Omega_2} \|\tilde k_s\|_{\infty} \leq
  d^{-1-s} \|f\|_{1,\Omega_2}.
  \end{equation*}
 Hence,
 \begin{align*}
  \|\lap^{\frac{s}{2}} f\|_{(p_1,q_1), \Omega_1} 
  & \leq |\Omega_1|^{1/p_1} \|\lap^{\frac{s}{2}} f\|_{\infty, \Omega_1}
  \leq d^{-1-s}|\Omega_1|^{1/p_1}\|f\|_{1, \Omega_2} \\ 
   & \leq d^{-1-s} |\Omega_1|^{1/p_1} |\Omega_2|^{1-1/p_2} \|f\|_{(p_2,q_2), \Omega_2}. \qedhere   
 \end{align*}
\end{proof}

A quite immediate consequence of this quasi-locality and Sobolev 
imbeddings is the following lemma. To state it, let
\begin{equation}\label{eq:annulus}
 A^k_{\Lambda,r}:= B_{2^{k}\Lambda r}- B_{2^{k-1} \Lambda r}.
\end{equation}

\begin{lemma}\label{lem:QuasiLocalityOfNorms}
Let $p \in (1,\infty), q,\tilde q \in [1,\infty]$, $s \in (-1,1/p)$ and $- \frac 1 {p^\ast} = s - \frac 1 p$.
\begin{enumerate}
  \item For $f \in\mathcal S(\R)$ with $\supp f \subset B_r$ we have
    \begin{equation*}
  \|\lap^{-\frac{s}{2}}f\|_{(p^\ast,q);\mathbb R - B_{\Lambda r}} 
  \aleq \Lambda^{-1+\fracm p} \|f\|_{(p, \tilde q),B_{r}}
 \end{equation*}
 uniformly for all $\Lambda >2$.
 
 \item If $s \in[0,\frac1p)$, we have for all $f \in\mathcal S(\R)$
 \begin{equation*}
  \|\lapms{s}f\|_{(p^\ast,q);\mathbb R - B_{\Lambda r}} 
  \aleq \Lambda^{-1+\fracm p} \|f\|_{(p, \tilde q);B_r} + \|f\|_{(p, q);\mathbb R - B_{r}}
 \end{equation*} 
 uniformly for all $\Lambda >2$.

 \item For $s\in[0,\frac1p)$ and any $f\in\mathcal S(\R)$, we have
 \begin{equation*}
    \|\lapms{s}f\|_{(p^\ast,q),B_r} 
    \aleq \|f\|_{(p,q),B_{\Lambda r}} + \Lambda^{-\frac 1{p^\ast}}
	\sum_{k=1}^\infty \left(2^{-k}\right)^{\frac 1 {p^\ast}}
	\|f\|_{(p,\tilde q),A^k_{\Lambda, r}}
 \end{equation*}
 uniformly for all $\Lambda >2$.
\end{enumerate}

% Especially,
% \[ \|\lapms{s}f\|_{(p^\ast,q),B_r} 
%  \aleq \sum_{k=1}^\infty \left(\Lambda^{1-k}\right)^{\frac 1 {p^\ast}}
%	\|f\|_{(p,q),B_{\Lambda^k r}}. \]
\end{lemma}

% \begin{lemma} \label{lem:QuasiLocalityOfNorms}
%  Let $p \in (1,\infty), q \in [1,\infty]$, $s \in [0,1/p)$ and $- \frac 1 {p^\ast} = s - \frac 1 p$.
%  For any $f\in\mathcal S(\R)$, we have
%  \begin{equation*}
%     \|\lapms{s}f\|_{(p^\ast,q),B_r} 
%     \aleq \|f\|_{(p,q),B_{\Lambda r}} + \Lambda^{-\frac 1{p^\ast}}
% 	\sum_{k=1}^\infty \left(2^{-k}\right)^{\frac 1 {p^\ast}}
% 	\|f\|_{(p,q),A^k_{\Lambda, r}}
%  \end{equation*}
%  uniformly for all $\Lambda >2$.
% % Especially,
% % \[ \|\lapms{s}f\|_{(p^\ast,q),B_r} 
% %  \aleq \sum_{k=1}^\infty \left(\Lambda^{1-k}\right)^{\frac 1 {p^\ast}}
% %	\|f\|_{(p,q),B_{\Lambda^k r}}. \]
% \end{lemma}

\begin{proof}
 For the first inequality we use Lemma~\ref{lem:QuasiLocality}
 and sum up the estimate
 \begin{equation*}
  \|\lap^{-\frac{s}{2}} f\|_{(p^*,  q ),A^k_{\Lambda, r}} \aleq (\Lambda 2^k)^{-1+\fracm p} \|f\|_{(p, \tilde q);B_r}.
 \end{equation*}

 In the second case, we use 
 \begin{align*}
  \|I_sf\|_{(p^\ast, q), \mathbb R - B_{\Lambda r}}
  &\leq \|I_s(\chi_{B_r}f)\|_{(p^\ast, q), \mathbb R - B_{\Lambda r}}
  + \|I_s((1-\chi_{B_r})f)\|_{(p^\ast, q), \mathbb R - B_{\Lambda r}}
% \\ 
% &\leq \sum_{k=1}^\infty \|I_s(\chi_{B_r}f)\|_{(p^\ast, q), A^k_{\Lambda r}}
%  + \|I_s((1-\chi_{B_r})f)\|_{(p^\ast, q), \mathbb R - B_{\Lambda r}}.
 \end{align*}
 and estimate the first term using~(i)
 % Lemma~\ref{lem:QuasiLocality}
 and the second term
 using Sobolev's inequality to get
 \begin{equation*}
  \|I_sf\|_{(p^\ast, q), \mathbb R - B_{\Lambda r}} 
%  \aleq \sum_{k=1}^\infty (\Lambda2^k)^{-1+\frac 1 p} \|f\|_{(p,\tilde q), B_r} 
%  + \|f\|_{(p,q); \mathbb R - B_{r}}
  \aleq \Lambda^{-1+\frac 1 p} \|f\|_{(p,\tilde q), B_r} 
  + \|f\|_{(p,q); \mathbb R - B_{r}}.
 \end{equation*}

 To deduce the last inequality, we decompose $f= \chi_{B_{\Lambda r}} f + \sum_{k=1}^\infty \chi_{A^{k}_{\Lambda, r}} f$ 
 and estimate using Sobolev's inequality
 \begin{align*}
  \|\lapms{s}(\chi_{B_{\Lambda r}}f)\|_{(p^\ast,q),B_r} \aleq \|f\|_{(p,q),B_{\Lambda r}}
 \end{align*}
 and using Lemma~\ref{lem:QuasiLocality}
 \begin{align*}
  \|\lapms{s}(\chi_{A^{k}_{\Lambda,r}}f)\|_{(p^\ast,q), B_r} 
  \aleq (\Lambda 2^k)^{-\frac 1 {p^\ast}}\|f\|_{(p,\tilde q),A^k_{\Lambda,r}}.
 \end{align*}
 Summing up, this proves the last inequality.
\end{proof}


% \begin{proof}
%  Since the inequality trivially holds for $s=0$ let us assume that $s>0$.
%  We decompose $f= \chi_{B_{\Lambda r}} f + \sum_{k=1}^\infty \chi_{A^{k}_{\Lambda, r}} f$ and estimate 
%  using Sobolev's inequality
%  \begin{align*}
%   \|\lapms{s}(\chi_{B_{\Lambda r}}f)\|_{(p^\ast,q),B_r} \aleq \|f\|_{(p,q),B_{\Lambda r}}
%  \end{align*}
%  and using Lemma~\ref{lem:QuasiLocality}
%  \begin{align*}
%   \|\lapms{s}(\chi_{A^{k}_{\Lambda,r}}f)\|_{(p^\ast,q), B_r} 
%   \aleq (\Lambda 2^k)^{-\frac 1 {p^\ast}}\|f\|_{(p,q),A^k_{\Lambda,r}}
%  \end{align*}
%  which finishes the proof.
% \end{proof}


Finally, we use the quasi-locality to prove

\begin{proposition}\label{pr:localcontrolelliptic}
For $p \in (1,\infty)$, $q \in [1,\infty]$, $s,t \geq 0$ with $0< s+t < 1$
there is a $\theta >0$ such that
we have for any $f \in L^{p,q}(\R)$, $\Lambda > 2$, and $r>0$  
\begin{multline*}
 \vrac{\laps{s} f}_{(p,q), B_r}  \aleq 
  \sup_{\ontop{\varphi \in C_0^\infty(B_{\Lambda^2 r}),}{\|\laps{t}\varphi\|_{(p',q')} \leq 1}}  \int f \laps{s+t} \varphi 
  +  r^{-s}\Lambda^{-\theta}\ \|f\|_{(p,q);B_{\Lambda r}} + r^{-s} 
   \Lambda^{-\theta} \sum_{l=1}^\infty 2^{-\theta l} \vrac{f}_{(p,q),A_{\Lambda, r}^l}.
\end{multline*}
\end{proposition}

\begin{proof}
Assume that 
\[
 \sup_{\varphi \in C_0^\infty(B_{\Lambda^{2} r}), \|\laps{t}\varphi\|_{(p',q')} \leq 1}
  \int f \laps{s+t} \varphi  \leq K.
\]
For $g \in C_0^{\infty}(B_r)$ we consider
\begin{align*}
 \int_{\mathbb R} (\laps{s} f) g dx= \int_{\mathbb R} f\laps{s} g dx
 = \int_{B_{\Lambda r}} f\laps{s} g dx+ \sum_{l=1}^\infty \int_{A^l_{\Lambda, r}} f\laps{s} g dx.
\end{align*}
For the first term we use a smooth partition of unity $\br{\eta_k}_{{k\in\N\cup\set0}}$
with $\supp \eta_0 \subset B_{\Lambda^2 r}$ and 
$\supp \eta_k \subset B_{(\Lambda^{2} 2^{k}+1)r} -B_{(\Lambda^{2} 2^{k-1}-1)r}$
to get
\begin{align*}
   \left|\int_{B_{\Lambda r}} f\laps{s} g dx\right| 
   &\leq \left|\int_{B_{\Lambda r}} f\laps{t+s} \br{\eta_0
  \lapms{t} g} dx\right|+ \sum_{l=1}^\infty
  \left|\int_{B_{\Lambda r}} f\laps{t+s} \br{\eta_l
  \lapms{t} g} dx\right| \\
  &\aleq K \| g\|_{(p',q')} + \|f\|_{(p,q); B_{\Lambda r}} \sum_{l=1}^\infty 
   \|\laps{t+s} \br{\eta_l
  \lapms{t} g}\|_{(p',q');B_{\Lambda r}}.
\end{align*}
Since by Lemma~\ref{lem:QuasiLocality} we have
\begin{align*}
 \|\laps{t+s} \eta_l \lapms{t} g\|_{(p',q');B_{\Lambda r}} 
 &\aleq (\Lambda^2 2^l r)^{-1-t-s} (\Lambda r)^{1/p'} (\Lambda^2 2^l r)^{1-1/p'}
 \|\eta_l\lapms{t} g\|_{(p',q')}  \\
 &\leq  (2^l )^{-t-s-1/p'} \Lambda^{-(2s+2t+1/p')} r^{-s-t}
 \|\lapms{t}g\|_{(p',q')}  
\end{align*}
we can estimate this further by
%\begin{align*}
\[  K \|g\|_{(p',q')} + \Lambda^{-\theta_{1}}r^{-s-t}\|f\|_{(p,q); B_{\Lambda r}} \|\lapms{t} g\|_{(p',q')} \]
% &=   K \|g\|_{(p',q')} + \Lambda^{-\theta_1}r^{-s-t}\|f\|_{(p,q); B_r}  \|\lapms{t} g\|_{(p',q')}.
%\end{align*}
for $\theta_1 = 2t+2s+1/p'$.

Since $\supp g \subset B_r$ we get a scaled Poincar\'e inequality by applying first the Sobolev and then the H\"older inequality
\begin{equation*}
 \|\lapms{t} g\|_{(p',q')} \aleq \norm{g}_{(\frac{p'}{p't+1},q'),B_r}
 \aleq  r^{t} \|g\|_{(p',q')}  
\end{equation*}
and using the quasi-locality (Lemma~\ref{lem:QuasiLocality}) once more
\begin{equation*}
 \|\laps{s} g\|_{(p',q'),A^k_{\Lambda r}}
  \aleq (2^{k}\Lambda r)^{-1-s}r^{1-1/p'}(2^{k}\Lambda r)^{1/p'}\|g\|_{(p',q')}
  = \Lambda^{-\theta_2} 2^{-k \theta_2} r^{-s} \|g\|_{(p',q')}
\end{equation*}
for $\theta_2 = 1+s-1/p'$.

Hence, 
\begin{equation*}
 \int_{\mathbb R} \br{\laps{s} f}\ g dx \leq \left(K + \Lambda ^{-\theta} r^{-s} 
 \|f\|_{(p,q);B_{\Lambda r}} + r^{-s} \Lambda^{-\theta}
  \sum_{l=1}^\infty 2^{-\theta l} 
  \|f\|_{(p,q);A^l_{\Lambda, r}} \right) \|g\|_{(p',q')},
\end{equation*}
for $\theta = \min\{\theta_1, \theta_2\}$ which by duality proves the proposition.
\end{proof}


\subsection{Proofs of Lemmata~\ref{la:tangentialpart} and~\ref{lem:CriticalTerm}}

The following lemma is the starting point for the 
estimates of $H_s$ and essentially follows from the mean value theorem or a first-order Taylor expansion.

\begin{lemma} \label{lem:MultiplierEstimate}
Let $\delta \in [0,1]$, $\alpha \in (0,1)$. Then for almost all $x,y, \xi \in \mathbb R^n$
we have
\begin{equation*}
 \left||x-\xi|^{-1+\alpha} - |y-\xi|^{-1+\alpha} \right| \aleq |x-y|^\delta 
 \left( |y-\xi|^{-1+\alpha-\delta} + |x-\xi|^{-1+\alpha-\delta} \chi_{|x-y|>2|x-\xi|}\right).
\end{equation*}
% \begin{equation*}
%  \left||x-\xi|^{-1+1/2} - |y-\xi|^{-1+1/2} \right| \aleq |x-y|^\delta 
%  \left( |y-\xi|^{-1+1/2-\delta} + |x-\xi|^{-1+1/2-\delta} \chi_{|x-y|>2|x-\xi|}\right).
% \end{equation*}
\end{lemma}

\begin{proof}
 If $|x-y|>2|x-\xi|$ we get
 \begin{equation*}
  |y-\xi| \geq |y-x|-|x-\xi| % \geq \frac 1 2 |x-y|
  > |x-\xi|
 \end{equation*}
 and hence
 \begin{equation*}
  \left||x-\xi|^{-1+\alpha} - |y-\xi|^{-1+\alpha} \right| \aleq |x-\xi|^{-1+\alpha}
  \aleq |x-y|^{\delta} |x-\xi|^{-1+\alpha-\delta}.
 \end{equation*}

 If $|x-y|\leq 2 |x-\xi|$ we first observe that the above argument leads to
 \begin{equation*}
   \left||x-\xi|^{-1+\alpha} - |y-\xi|^{-1+\alpha} \right|
  \aleq |x-y|^{\delta} |y-\xi|^{-1+\alpha-\delta}
 \end{equation*}
 if $|x-y|>2|y-\xi|$.

 To deal with the case that both $|x-y|\leq 2 |x-\xi|$ and $|x-y|\leq 2 |y-\xi|$
 we observe that then
 \begin{equation*}
  |y-\xi| \leq |x-\xi| + |x-y| \leq 3 |x-\xi|.
 \end{equation*}
 Hence, we get using the mean value theorem
 \begin{align*}
   \left||x-\xi|^{-1+\alpha} - |y-\xi|^{-1+\alpha} \right|
   &\aleq |x-y| \max\left\{|x-\xi|^{-2+\alpha}, |y-\xi|^{-2+\alpha}\right\}  \\
   &\aleq |x-y| |y-\xi|^{-2+\alpha} \aleq |x-y|^\delta |y-\xi|^{-1+\alpha-\delta}.\qedhere
 \end{align*}
\end{proof}


We use the lemma above to derive the following pointwise estimate for $H_{1/2+s}$.

\begin{lemma}[\cite{SchikorraIntBoundFrac11}]\label{lem:EstimateForH12s}
For $s \in [0,\frac12)$ %there exists $\varepsilon_0 >0 $ such that for
and functions $a,b\in\mathcal S(\R)$
% \begin{equation}\label{eq:ab}
%  a=\lapmv  \lapv  a \qquad\text{and}\qquad b= \lapmv  \lapv  b,
% \end{equation}
the following holds for any $\varepsilon,\eps' \in [0,\frac16-\frac s3)$:
\begin{align*}
 |H_{1/2+s}(a,b)| &\aleq \lapms{1/6-s/3} (\lapms{1/6-s/3-\varepsilon}|\laps{1/2-\varepsilon} a| \; \lapms{1/6-s/3-\varepsilon'}|\laps{1/2-\varepsilon'} b|) 
   \\ &\quad+ \lapms{1/3-2s/3-\varepsilon}|\laps{1/2-\varepsilon} a| \; \lapms{1/6-s/3-\varepsilon'}|\laps{1/2-\varepsilon'} b|
   \\ & \quad+\lapms{1/6-s/3-\varepsilon}|\laps{1/2-\varepsilon} a|\; \lapms{1/3-2s/3-\varepsilon'}|\laps{1/2-\varepsilon'} b|
   \\ & \quad+\lapms{1/4-s/2-\varepsilon}|\laps{1/2-\varepsilon} a| \; \lapms{1/4-s/2-\varepsilon'}|\laps{1/2-\varepsilon'} b|. 
\end{align*}
\end{lemma}

\begin{proof}
%Set $\varepsilon_0 := \fracm{6} - \frac{s}{3}$, and let $\varepsilon \in [0,\varepsilon_0)$, then the following arguments holds.
In order to shorten notation, we restrict to $\eps'=\eps$;
the general case is parallel.
 We use the identities $a=\lapms{\fracm{2}-\varepsilon}  \laps{\fracm{2}-\varepsilon}a$ and $b = \lapms{\fracm{2}-\varepsilon}  \laps{\fracm{2}-\varepsilon} b$.
 %For simplicity assume now $\eps = 0$.
 Then,
\begin{align*}
 &|H_{1/2+s}(a,b)(x)| \\&\aleq \int \int \int \frac {\left(|y-z|^{-1+1/2-\eps}-|x-z|^{-1+1/2-\eps} \right) 
  \left( |y-w|^{-1+1/2-\eps} - |x-w|^{-1+1/2-\eps} \right)}{|y-x|^{1+ (1/2+s)}}\cdot{}\\
  &\qquad\qquad\qquad{}\cdot\abs{\laps{\frac{1-2\eps}2} a(z)} \,\abs{\laps{\frac{1-2\eps}2} b(w)} \d w \d z \d y .
\end{align*}
Applying Lemma~\ref{lem:MultiplierEstimate} we get
\begin{align*}
 & \frac {\left(|y-z|^{-1+1/2-\eps}-|x-z|^{-1+1/2-\eps} \right) 
  \left( |y-w|^{-1+1/2-\eps} - |x-w|^{-1+1/2-\eps} \right)}{|y-x|^{1+ (1/2+s)}} \\
 &\aleq \frac {\scriptstyle\left(|y-w|^{-1+1/2-\eps-\delta} + |x-w|^{-1+1/2-\eps-\delta} \chi_{|x-y|>2|x-w|} \right) 
		\left(|y-z|^{-1+1/2-\eps-\delta} + |x-z|^{-1+1/2-\eps-\delta} \chi_{|x-y|>2|x-z|} \right)}
	     {|x-y|^{1+1/2+s-2\delta}} \\
 &\aleq \frac {(|y-w||y-z|)^{-1+1/2-\eps-\delta} + (|x-w||y-z|)^{-1+1/2-\eps-\delta}  
  + (|y-w||x-z|)^{-1+1/2-\eps-\delta}}{|x-y|^{1+1/2+s-2\delta}} \\
 & \quad{} + \frac {(|x-w||x-z|)^{-1+1/2-\eps-\delta}}{|x-y|^{1+1/2+s-2\delta}}
    \chi_{|x-y|>2|x-w|} \ \chi_{|x-y|>2|x-z|}.
\end{align*}
For $\delta =1/3 + s/3$ we hence get
\begin{align*}
 |H_{1/2+s}(a,b)|& \aleq \lapms{1/6-s/3} (\lapms{1/6-s/3-\eps}|\laps{\frac{1-2\eps}2} a| \ \lapms{1/6-s/3-\eps}|\laps{\frac{1-2\eps}2} b|) 
  \\ & \quad+\lapms{1/6-s/3-\eps}|\laps{\frac{1-2\eps}2} a|  \ \lapms{1/3-2s/3-\eps}|\laps{\frac{1-2\eps}2} b|
  \\ &\quad+ \lapms{1/3-2s/3-\eps}|\laps{\frac{1-2\eps}2} a| \ \lapms{1/6-s/3-\eps}|\laps{\frac{1-2\eps}2} b|
  + A
\end{align*}
where
\begin{align*}
 A &\aleq \int \int \brac{\int \frac {(|x-w||x-z|)^{-1+1/2-\eps-\delta}}{|x-y|^{1+1/2+s-2\delta}}
    \chi_{|x-y|>2|x-w|} \  \chi_{|x-y|>2|x-z|}\ \d y} \	 |\laps{\frac{1-2\eps}2}  a(z)|\ |\laps{\frac{1-2\eps}2}  b(w)| \d w \d z \\
   &\aleq \int \int   |x-w|^{-1+1/4-\eps-s/2 } |x-z|^{-1+1/4-\eps-s/2} |\laps{\frac{1-2\eps}2}  a(z)|\ |\laps{\frac{1-2\eps}2}  b(w)| \d w \d z \\
   &\aleq \lapms{1/4-s/2-\eps}|\laps{\frac{1-2\eps}2} a|\ \lapms{1/4-s/2-\eps}| \laps{\frac{1-2\eps}2} b|. \qedhere
\end{align*}
\end{proof}

To estimate this further, we will use the following fact about {\em lower order products}
which we get using the quasi-locality.

% \begin{lemma} [Lower order products] \label{lem:LowerOrderProducts}
%  Let $0 \leq s_{1,\nu}, s_{2,\nu}, s_{3,\nu} < 1/2$ with
%  $s_{1,\nu}+s_{2,\nu}+s_{3,\nu} = 1/2-s$ for some $s\in [0,1/2)$, $s + s_1 < \fracm{2}$, $q\in[1,\infty]$.
%  There exists some $\theta > 0$ such that, if for any bilinear operator $G$ and $u,v\in\mathcal S(\R)$ the following pointwise estimate holds,
%  \begin{equation*}
%   |G(u,v)| \aleq \sum_{\nu=1}^N \lapms{s_{1,\nu}} \left( \lapms{s_{2,\nu}}  |u|\ 
%   \lapms{s_{3,\nu}} |v|\right),
%  \end{equation*}
%  then we have for any $\Lambda>2$, $\varepsilon \in [0,s_3]$, $\delta \in [0,s_2]$, and $\fracm{q} = \fracm{q_1}+ \fracm{q_2}$
%  \begin{align}\label{eq:Guvest}
%     &\|G(u,v)\|_{(\frac 2 {1+2s}, q), B_r} \\ \nonumber
%   &\aleq \|\lapms{\delta}\abs{u}\|_{(\frac2{1-2\delta},q_1),B_{\Lambda r}} \|\lapms{\eps}\abs{v}\|_{(\frac2{1-2\eps},q_2),B_{\Lambda r}} + \|\lapms{\delta}\abs{u}\|_{(\frac2{1-2\delta},q_1)}
%     \sum_{k=2}^\infty \Lambda^{-\theta(k-1)}\|\lapms{\eps}\abs{v}\|_{(\frac2{1-2\eps},q_2),B_{\Lambda^k r}}.
%  \end{align}
% If $\supp \lapmv v \subset B_r$ we furthermore get for any $k \geq 2$, $\Lambda \geq 16$,
%  \begin{equation}\label{eq:Ginversedssupport}
%   \|G(u,v)\|_{(\frac2{1+2s},q),A^k_{\Lambda,r}} \aleq \Lambda^{-\theta k} \vrac{u}_{(2,\infty),\R}\ \|\lapms{\eps} \abs{v}\|_{(\frac 2 {1-2\eps}, 1),\R}.
%  \end{equation}
% \end{lemma}
% 
% \begin{proof}
%  It is enough to prove the statements for the case $N=1$ and we set in this case
%  $s_i:=s_{i,1}$ for $i=1,2,3$. Note that in case $s=0$ at least two of the parameters $s_1$, $s_2$, $s_3$ are positive.
%  
%  Now we use several times the quasi-locality (Lemma~\ref{lem:QuasiLocalityOfNorms}, e.g., in the first application with $\Lambda^{1/2}$ in place of $\Lambda$, for $p^*\hat=\frac{2}{1+2s}$, $q\hat=\infty$, $s\hat=s_1$, $p\hat=\frac{2}{1+2(s+s_1)}$), H\"older's and Sobolev's inequalities. Let $\eps \in [0,s_3]$, $\delta \in [0,s_2]$,
%  \begin{align*}
%   &\|G(u,v)\|_{(\frac 2 {1+2s},q),B_r} \\
%   &\aleq \sum_{k=1}^\infty (\Lambda ^{(k-1)/2})^{-1/2 - s} 
%   \|\lapms{s_2}\abs{u}\ \lapms{s_3}\abs{v}\|_{(\frac 2{1+2(s+s_1)},q), B_{\Lambda^{k/2} r}}\\
%   &\aleq \sum_{k=1}^\infty (\Lambda ^{(k-1)/2})^{-1/2 - s} 
% 	\|\lapms{s_2}\abs{u}\|_{(\frac 2 {1-{2s_2}},q_1),B_{\Lambda^{k/2} r}}\  
% 	\|\lapms{s_3}\abs{v}\|_{(\frac 2 {1-2s_3}, q_2), B_{\Lambda^{k/2} r}}
%   \\ 
%   &\aleq
% 	\sum_{k=1}^\infty \sum_{l_1=1}^\infty \sum_{l_2 =1}^\infty(\Lambda ^{(k-1)/2})^{-1/2 - s}
% 	(\Lambda ^{(l_1-1)/2})^{-1/2 + s_2} (\Lambda^{(l_2 -1)/2})^{-1/2 + s_3} \cdot\\
% 	&\qquad\qquad\qquad{}\cdot\|\lapms{\delta}\abs{u}\|_{(\frac{2}{1-2\delta}, q_1 ), B_{\Lambda^{(k+l_1)/2} r}}
% 	\|\lapms{\eps} \abs{v}\|_{(\frac2{1-2\eps}, q_2), B_{\Lambda^{(k+l_2)/2} r}}
%   \\
%   &\aleq \|\lapms{\delta}\abs{u}\|_{(\frac{2}{1-2\delta},q_1),B_{\Lambda r}}  \|\lapms{\eps} \abs{v}\|_{(\frac2{1-2\eps},q_2),B_{\Lambda r}}
%     + \|\lapms{\varepsilon}\abs{u}\|_{(\frac{2}{1-2\delta}.q_1)} 
% 	\sum_{k=2}^\infty (\Lambda^{k-1})^{-\theta}\| \lapms{\varepsilon} \abs{v}\|_{(\frac2{1-2\varepsilon},q_2),B_{\Lambda^{k}r}},
%  \end{align*}
% for $\theta >0$ small enough, depending on $s_1$, $s_2$, $s_3$, but not on $\varepsilon$. 
% As for the second estimate, using Lemma~\ref{lem:QuasiLocality} and again Lemma~\ref{lem:QuasiLocalityOfNorms}, we have for $k \geq 1$,
% \begin{align*}
%     \|G(u,v)\|_{(\frac2{1+2s},q),A^k_{\Lambda,r}} 
%     &= \|\lapms{s_1}\left( \lapms{s_2} |u| \ \lapms{s_3} |v| \right)\|_{(\frac2{1+2s},q),A^k_{\Lambda,r}}
%   \\
%     &\aleq \Lambda^{-k \frac{1/2-s_1-s}{2}} \|\lapms{s_2} |u| \ \lapms{s_3} |v|\|_{(\frac 2 {1+2(s+s_1)}, \infty), B_{\Lambda^{k/2} r}}
%     \\
%     &\quad+ \|\lapms{s_2} |u| \ \lapms{s_3} |v|\|_{(\frac 2 {1+2(s+s_1)},q), B_{\Lambda^{k+2}r}-B_{\Lambda^{k/2}r}}
%     \\
%     &\quad + \sum_{l=2}^\infty \|\lapms{s_2} |u| \ \lapms{s_3} |v|\|_{(\frac 2 {1+2(s+s_1)}, q), A^{k+l}_{\Lambda,r}}
%   \\
%     &\aleq \Lambda^{-k \frac{1/2-s_1-s}{2}} \|\lapms{s_2} |u| \|_{(\frac 2 {1-2s_2},\infty),\R}\  
% 	  \|\lapms{s_3} |v|\|_{(\frac 2 {1-2s_3}, \infty),\R}
%     \\
%     &\quad+ \|\lapms{s_2} |u|\|_{(\frac 2 {1-2 s_2}, \infty),\R} \ 
%     \| \lapms{s_3} |v|\|_{(\frac 2 {1-2s_3},q), B_{\Lambda^{k+2}r}-B_{\Lambda^{k/2}r}}
%     \\
%     &\quad + \sum_{l=2}^\infty \|\lapms{s_2} |u| \|_{(\frac 2 {1-2s_2},\infty),\R}\ 
%       \| \lapms{s_3} |v|\|_{(\frac 2 {1-2s_3}, q), A^{k+l}_{\Lambda,r}}
%   \\
%   %
%   %
%   &\aleq \Lambda^{-k \frac{1/2-s_1-s}{2}} \|u\|_{(2,\infty)} \|\lapms{s_3} v\|_{(\frac 2 {1-2s_3}, \infty),\R} \\
%   &\quad+ \| u\|_{(2, \infty),\R} \ 
%     \| \lapms{s_3} v\|_{(\frac 2 {1-2s_3},q), B_{\Lambda^{k+2}r}-B_{\Lambda^{k/2}r}}
%     \\
%     &\quad + \sum_{l=2}^\infty \| u \|_{(2,\infty),\R}\ 
%       \| \lapms{s_3} v\|_{(\frac 2 {1-2s_3}, q), A^{k+l}_{\Lambda,r}}.
% %     
% %     
% %     
% %     & \quad + \Lambda^{-k(1/2-s_3)}\|u\|_{(2,\infty)} \|v\|_{2,\infty}
% %      + \sum_{l=2}^\infty \Lambda^{-(k+l)(1/2-s_3)} \|u\|_{(2,\infty)} \|v\|_{(2,\infty)}.
%  \end{align*}
%  Since
%  \[
%   \lapms{s_3} v = \laps{\frac{1}{2}-s_3} \lapmv v,
%  \]
% under the assumption $\supp \lapmv v \subset \overline{B_r}$, Lemma~\ref{lem:QuasiLocality} with $s \hat= \fracm{2} - s_3$, $p_1 \hat= \frac{2}{1-2s_3}$, $p_2 \hat= \infty$, $\Omega_3 \hat= B_r$ implies (note this is a rather rough, and non-optimal Lorentz-space estimate, which nonetheless suffices for our purposes!) 
% \[
%  \| \lapms{s_3} v\|_{(\frac 2 {1-2s_3},q), B_{\Lambda^{k+2}r}-B_{\Lambda^{k/2}r}} \aleq \Lambda^{-k(\fracm{4}+\frac{s_3}{2})+\fracm{2}-s_3}\ \vrac{\lapmv v}_{\infty} \aleq \Lambda^{-k(\fracm{4}+\frac{s_3}{2})+\fracm{2}-s_3}\ \vrac{\lapms{s_3} v}_{(\frac 2 {1-2s_3},1)},
% \]
% and
% \[
% \| \lapms{s_3} v\|_{(\frac 2 {1-2s_3}, q), A^{k+l}_{\Lambda,r}} \aleq \Lambda^{-(k+l)+\fracm{2}-s_3} \vrac{\lapms{s_3} v}_{(\frac 2 {1-2s_3},1)}.
% \]
% We thus have shown,
% \begin{align*}
%     \|G(u,v)\|_{(\frac2{1+2s},q),A^k_{\Lambda,r}} 
%     &\aleq \brac{\Lambda^{-k \frac{1/2-s_1-s}{2}} +  \Lambda^{-k(\fracm{4}+\frac{s_3}{2})+\fracm{2}-s_3} + \Lambda^{-k+\fracm{2}-s_3}}\ \|u\|_{(2,\infty)} \|\lapms{s_3} v\|_{(\frac 2 {1-2s_3}, 1),\R} \\
% \end{align*}
% This implies \eqref{eq:Ginversedssupport} for a $\theta > 0$ depending on $s_1$,$s_2$,$s_3$, but not on $\varepsilon \in [0,s_3]$.
% \end{proof}
\begin{lemma} [Lower order products] \label{lem:LowerOrderProducts.v2}
  Let for $s > 0$, $0 \leq s_{1}, s_{2}, s_{3} \leq 1/2$, $s_{1} + s_{2}+ s_{3} = s$, and at least two of these $s_i$, $i=1,2,3$, non-zero. Let $p \in (1,\infty)$, $p_2 \in (1,\fracm{s_2}),p_3 \in (1,\fracm{s_3})$ and such that
  \[
   \fracm{p} = \fracm{p_2} + \fracm{p_3} - s.
  \]
Assume moreover that
\begin{equation*}
  |G(u,v)| \aleq \lapms{s_{1}} \left( \lapms{s_{2}}  |u|\ 
  \lapms{s_{3}} |v|\right),
 \end{equation*}
 then there is some $\theta>0$ such that
 we have for any $\Lambda>4$, and $\fracm{q} = \fracm{q_1}+ \fracm{q_2}$,
 \begin{align}\label{eq:Guvest}
    &\|G(u,v)\|_{(p, q), B_r} \\ \nonumber
  &\aleq \|u\|_{(p_2,q_1),B_{\Lambda r}} \|v\|_{(p_3,q_2),B_{\Lambda r}} + \|u\|_{(p_2,q_1)}
    \Lambda^{-\theta}\sum_{k=2}^\infty 2^{-\theta k}\|v\|_{(p_3,q_2),B_{\Lambda 2^k r}}.
 \end{align}
 If $\supp \laps{t}v  \subset \overline{B_r}$ for some $t \in [0,\fracm{2}]$, 
 we furthermore get for any $k \geq 2$, $\Lambda > 16$,
 \begin{equation}\label{eq:Ginversedssupport}
  \|G(u,v)\|_{(p,q),A^k_{\Lambda,r}} 
  \aleq \Lambda^{-\theta} 2^{-\theta k} \vrac{u}_{(p_2,\infty),\R}\ \|v\|_{(p_3, 1),\R}.
 \end{equation}
\end{lemma}
\begin{proof}
Let
\[
  \fracm{p_2^\ast} :=  \fracm{p_2} -s_2, 
\]
\[
  \fracm{p_3^\ast} :=  \fracm{p_3} -s_3, 
 \]
\[
  \fracm{p_1} :=  \fracm{p_2^\ast} + \fracm{p_3^\ast} = \fracm{p_2} + \fracm{p_3} - s_2 - s_3,
\]
and
\[
 \fracm{p_1^\ast} := \fracm{p} \equiv \fracm{p_2} + \fracm{p_3} - s
 =\frac{1}{p_{1}}-s_{1}.
\]

Then by Lemma~\ref{lem:QuasiLocalityOfNorms}~(iii), for $\Lambda \hat = \sqrt{\Lambda}$, $p^\ast \hat = p_1^\ast=p$, $p \hat= p_1$, $s\hat=s_{1}$
 \begin{align*}
  &\|G(u,v)\|_{(p_1^\ast,q),B_r} \aleq
  \vrac{\lapms{s_2}\abs{u} \ \lapms{s_3}\abs{v}}_{(p_1,q), B_{\Lambda^{1/2} r}} + \Lambda^{-1/2p}
  \sum_{k=1}^\infty 2 ^{-k/p} 
  \vrac{\lapms{s_2}\abs{u} \ \lapms{s_3}\abs{v}}_{(p_1,q), B_{\Lambda^{1/2}2^k r}} \\
  &\aleq
 \vrac{\lapms{s_2}\abs{u}}_{(p_2^\ast,q_1), B_{\Lambda^{1/2} r}} \ \vrac{\lapms{s_3}\abs{v}}_{(p_3^\ast,q_2), B_{\Lambda^{1/2} r}}+ \Lambda^{-1/2p}
 \sum_{k=1}^\infty 2 ^{-k/p} 
\vrac{\lapms{s_2}\abs{u}}_{(p_2^\ast,q_1), B_{2^k\Lambda^{1/2} r}} \ \vrac{\lapms{s_3}\abs{v}}_{(p_3^\ast,q_2), B_{2^k\Lambda^{1/2} r}}.
 \end{align*}
Applying yet again Lemma~\ref{lem:QuasiLocalityOfNorms}~(iii) to both factors in each sum with $\Lambda \hat = \sqrt{\Lambda}$
we arrive at
\begin{align*}
    \|G(u,v)\|_{(p, q), B_r} \aleq \|u\|_{(p_2,q_1),B_{\Lambda r}} \|v\|_{(p_3,q_2),B_{\Lambda r}} + \|u\|_{(p_2,q_1)}
    \Lambda^{-\theta}\sum_{k=2}^\infty 2^{-\theta k}\|v\|_{(p_3,q_2),B_{\Lambda 2^k r}}.
\end{align*}

As for the second estimate, we apply Lemma~\ref{lem:QuasiLocalityOfNorms}~(ii)
with $\Lambda \hat = \sqrt{\Lambda}$ to get
\begin{align*}
 \|G(u,v)\|_{(p,q), \mathbb R - B_{\Lambda r}} 
 &=\|\lapms{s_1}\left( \lapms{s_2} |u| \ \lapms{s_3} |v| \right)\|_{(p_1^\ast,q),\mathbb R - B_{\Lambda r}}
 \\
 &\aleq \Lambda^{-\theta_1} 
 \|\lapms{s_2} |u| \ \lapms{s_3} |v|\|_{(p_1, \infty), B_{\Lambda^{1/2} r}}
    + \|\lapms{s_2} |u| \ \lapms{s_3} |v|\|_{(p_1,q), \mathbb R -B_{\Lambda^{1/2}r}}
\end{align*}
where $\theta_{1} = \frac 1 2 - \fracm{2p_1}$.
We estimate the first term using H\"older's inequality and Sobolev imbedding theorem to get
\begin{equation*}
 \Lambda^{-\theta_1} 
 \|\lapms{s_2} |u| \ \lapms{s_3} |v|\|_{(p_1, \infty), B_{\Lambda^{1/2} r}}
 \leq \Lambda^{-\theta_1}\|u\|_{(p_2,\infty)} \|v\|_{(p_3,1)}.
\end{equation*}
For the second term we use H\"older's inequality and Sobolev's imbedding theorem to get
\begin{align*} 
 \|\lapms{s_2} |u| \ \lapms{s_3} |v|\|_{(p_1, q), \mathbb R - B_{\Lambda^{1/2} r}}
 &\aleq \|\lapms{s_2}\abs u\|_{(p^\ast_2,q_1), \mathbb R - B_{\Lambda^{1/2}r}} 
 \|\lapms{s_3}\abs v\|_{(p^\ast_3,q_2), \mathbb R - B_{\Lambda^{1/2}r}} \\
 &\aleq \|u\|_{(p_2,\infty), \mathbb R } 
 \|\lapms{s_3}|v|\|_{(p^\ast_3,q_{2}), \mathbb R - B_{\Lambda^{1/2}r}}.
\end{align*}
Since
\begin{align*}
 \|\lapms{s_3} |v|\|_{(p^\ast_3, q_2), \mathbb R - B_{\Lambda^{1/2} r}}
 & \aleq \Lambda^{-\theta_2} \|v\|_{(p_3,1);B_{\Lambda^{1/4}}} + \|v\|_{(p_3,q);\mathbb R - B_{\Lambda^{1/4}}}
\end{align*}
where $\theta_2 =  \fracm 4 - \fracm {4p_3}$ and,
using Lemma~\ref{lem:QuasiLocalityOfNorms}~(i)
\begin{equation*}
\|v\|_{(p_3,q);\mathbb R - B_{\Lambda^{1/4}}}
= \|\laps{t} (\lapms{t}v)\|_{(p_3,q);\mathbb R - B_{\Lambda^{1/4}}}
\leq \Lambda^{-\theta_2} \|I_t v\|_{(p^{\ast \ast}_3,1)} 
\leq \Lambda^{-\theta_2} \|I_t v\|_{(p_3,1)}
\end{equation*}
we deduce the statement for $\theta := \min\{\theta_1, \theta_2\}$.
% 
% As for the second estimate, using Lemma~\ref{lem:QuasiLocality} and
% proceeding as in the proof of Lemma~\ref{lem:QuasiLocalityOfNorms}, we have for $k \geq 1$, $\sigma \in (1,2)$, $\theta_{p_{1},\sigma}:=(1-\frac1{p_{1}})(1-\frac1\sigma)>0$
% \begin{align*}
%     \|G(u,v)\|_{(p,q),A^k_{\Lambda,r}} 
%     &= \|\lapms{s_1}\left( \lapms{s_2} |u| \ \lapms{s_3} |v| \right)\|_{(p_1^\ast,q),A^k_{\Lambda,r}}
%   \\
%     &\aleq \Lambda^{-\theta_{p_1,\sigma}k} \|\lapms{s_2} |u| \ \lapms{s_3} |v|\|_{(p_1, \infty), B_{\Lambda^{k/\sigma} r}}
%     \\
%     &\quad+ \|\lapms{s_2} |u| \ \lapms{s_3} |v|\|_{(p_1,q), B_{\Lambda^{k+2}r}-B_{\Lambda^{k/\sigma}r}}
%     \\
%     &\quad + \sum_{l=2}^\infty \|\lapms{s_2} |u| \ \lapms{s_3} |v|\|_{(p_1, q), A^{k+l}_{\Lambda,r}}
%   \\
%     &\aleq \Lambda^{-\theta_{p_1,\sigma}k} \|\lapms{s_2} |u| \|_{(p_2^\ast,\infty),\R}\  
% 	  \|\lapms{s_3} |v|\|_{(p_3^\ast, \infty),\R}
%     \\
%     &\quad+ \|\lapms{s_2} |u|\|_{(p_2^\ast, \infty),\R} \ 
%     \| \lapms{s_3} |v|\|_{(p_3^\ast,q), B_{\Lambda^{k+2}r}-B_{\Lambda^{k/\sigma}r}}
%     \\
%     &\quad + \sum_{l=2}^\infty \|\lapms{s_2} |u| \|_{(p_2^\ast,\infty),\R}\ 
%       \| \lapms{s_3} |v|\|_{(p_3^\ast, q), A^{k+l}_{\Lambda,r}}
%   \\
%   %
%   %
%   &\aleq \Lambda^{-\theta_{p_1,\sigma}k} \|u\|_{(p_2,\infty)}\ \|v\|_{(p_3, \infty),\R} \\
%   &\quad+ \| u\|_{(p_2, \infty),\R} \ 
%     \| \lapms{s_3} \abs{v}\|_{(p_3^\ast,q), B_{\Lambda^{k+2}r}-B_{\Lambda^{k/\sigma}r}}
%     \\
%     &\quad + \sum_{l=2}^\infty \| u \|_{(p_2,\infty),\R}\ 
%       \| \lapms{s_3} \abs{v}\|_{(p_3^\ast, q), A^{k+l}_{\Lambda,r}}.
% %     
% %     
% %     
% %     & \quad + \Lambda^{-k(1/2-s_3)}\|u\|_{(2,\infty)} \|v\|_{2,\infty}
% %      + \sum_{l=2}^\infty \Lambda^{-(k+l)(1/2-s_3)} \|u\|_{(2,\infty)} \|v\|_{(2,\infty)}.
%  \end{align*}
%  Note that $\theta_{p_{1},\sigma} \to 0$ as $\sigma \searrow 1$.  Since $\supp \lapms{t} v \subset \overline{B_r}$, we can write, for $\lambda > 1$,
%  \[
%   \lapms{s_3} \abs{v} = \lapms{s_3} \abs{\laps{t} \lapms{t} v} = \lapms{s_3} \abs{\laps{t} \brac{\chi_{B_{r}}\lapms{t} v}} 
%  \aleq \lapms{s_3} \abs{\chi_{B_{\lambda r}} v}  + \sum_{i=1}^\infty \lapms{s_3} \abs{\chi_{A^i_{\lambda r}}\laps{t} \brac{\chi_{B_{r}}\lapms{t} v}}.	 
%  \]
%  Using again Lemma~\ref{lem:QuasiLocality}, choosing $\sigma$ closer to $1$ as $p_3 \to 1$, we conclude.
% % 
% 
% 
% 
% 
% under the assumption $\supp \lapmv v \subset \overline{B_r}$, Lemma~\ref{lem:QuasiLocality} with $s \hat= \fracm{2} - s_3$, $p_1 \hat= \frac{2}{1-2s_3}$, $p_2 \hat= \infty$, $\Omega_3 \hat= B_r$ implies (note this is a rather rough, and non-optimal Lorentz-space estimate, which nonetheless suffices for our purposes!) 
% \[
%  \| \lapms{s_3} v\|_{(\frac 2 {1-2s_3},q), B_{\Lambda^{k+2}r}-B_{\Lambda^{k/2}r}} \aleq \Lambda^{-k(\fracm{4}+\frac{s_3}{2})+\fracm{2}-s_3}\ \vrac{\lapmv v}_{\infty} \aleq \Lambda^{-k(\fracm{4}+\frac{s_3}{2})+\fracm{2}-s_3}\ \vrac{\lapms{s_3} v}_{(\frac 2 {1-2s_3},1)},
% \]
% and
% \[
% \| \lapms{s_3} v\|_{(\frac 2 {1-2s_3}, q), A^{k+l}_{\Lambda,r}} \aleq \Lambda^{-(k+l)+\fracm{2}-s_3} \vrac{\lapms{s_3} v}_{(\frac 2 {1-2s_3},1)}.
% \]
% We thus have shown,
% \begin{align*}
%     \|G(u,v)\|_{(\frac2{1+2s},q),A^k_{\Lambda,r}} 
%     &\aleq \brac{\Lambda^{-k \frac{1/2-s_1-s}{2}} +  \Lambda^{-k(\fracm{4}+\frac{s_3}{2})+\fracm{2}-s_3} + \Lambda^{-k+\fracm{2}-s_3}}\ \|u\|_{(2,\infty)} \|\lapms{s_3} v\|_{(\frac 2 {1-2s_3}, 1),\R} \\
% \end{align*}
% This implies \eqref{eq:Ginversedssupport} for a $\theta > 0$ depending on $s_1$,$s_2$,$s_3$, but not on $\varepsilon \in [0,s_3]$.
\end{proof}


We then have the following
\begin{lemma}\label{lem:localizedHEst}
There is a $\varepsilon_{0} > 0$ such that
for $\delta,\eps\in[0,\varepsilon_{0})$,
for any $a,b\in\mathcal S(\R)$, $\Lambda>4$
\begin{align*}
    &\|H_{\frac{1}{2}}(a,b)\|_{(2, q), B_r} 
  \aleq \|\laps{1/2-\delta} a\|_{(\frac2{1-2\delta},q_1),B_{\Lambda r}} \|\laps{1/2-\eps} b\|_{(\frac2{1-2\eps},q_2),B_{\Lambda r}} + \\
  &\qquad\qquad\qquad\qquad{}+\|\laps{1/2-\delta} a\|_{(\frac2{1-2\delta},q_1)} \Lambda^{-\theta}
    \sum_{k=1}^\infty 2^{-\theta(k-1)}\|\laps{1/2-\eps} b\|_{(\frac2{1-2\eps},q_2),B_{2^k\Lambda r}}.
 \end{align*}
If $\supp b \subset B_r$ we furthermore get for any $k \geq 2$, $\Lambda > 16$, $s\in[0,\frac12)$
 \begin{equation*}
  \|H_{\fracm{2}}(a,b)\|_{(\frac2{1+2s},q),A^k_{\Lambda,r}} \aleq (\Lambda 2^k)^{-\theta} \vrac{\lapv a}_{(2,\infty),\R}\ \|\laps{1/2-\eps} b \|_{(\frac 2 {1-2\eps}, 1),\R}.
 \end{equation*}
 \end{lemma}
 
\begin{proof}
Immediately from Lemma~\ref{lem:EstimateForH12s}, Lemma~\ref{lem:LowerOrderProducts.v2} where $s = \fracm{2} - \varepsilon - \delta$. 
\end{proof}


We use the lemma above to estimate 
the normal part as stated in Lemma~\ref{la:tangentialpart}. 

\begin{proof}[Proof of Lemma \ref{la:tangentialpart}]
With~\eqref{eq:Dg2} we obtain $\norm{\laps{s}\sabs{g'}^2}_{(\frac2{1+2s},\infty),B_{r}}\aleq r^{\frac12+s}$.
In order to show~\eqref{eq:orthogpartest} by the decomposition~\eqref{eq:orthogpartHdec} it remains to treat the $H$-term.
We rewrite
\[
\begin{ma}
 \laps{s} H_{\frac{1}{2}}(a,b) 
&=& \underbrace{H_{s+\fracm{2}}(a,b)}_{=:I}
+ \underbrace{a\lapsv{1+2s}b - \laps{s}(a \lapv b)}_{=:II}
+ \underbrace{b\lapsv{1+2s}a - \laps{s}(b \lapv a)}_{=:III}.
\end{ma}
\]
We will show that all three terms satisfy the hypotheses of Lemma~\ref{lem:LowerOrderProducts.v2}.
Due to Lemma~\ref{lem:EstimateForH12s} this is true for the term $I$.

As for $II$ note that $II=0$ in case $s=0$. If $s\in(0,1)$,
the potential definition~\eqref{eq:fraclap-} of $\laps{s}$ gives
\begin{align*}
 {II(x)} &= c\br{a(x)\int \frac{\lapv b(y)-\lapv b(x)}{\abs{x-y}^{1+s}}\ \d y - \int \frac{a(y)\lapv b(y)-a(x)\lapv b(x)}{\abs{x-y}^{1+s}}\ \d y} \\
 &= c\int \frac{\br{a(x)-a(y)}\lapv b(y)}{\abs{x-y}^{1+s}}\ \d y.
\end{align*}
Using $a = \lapmv \lapv a$ we arrive at
\[
 \abs{II(x)} \aleq \int \frac{\abs{a(y)-a(x)}\abs{\lapv b(y)}}{\abs{x-y}^{1+s}}\ \d y \aleq \int \int \frac{\abs{\abs{y-z}^{-1+\frac{1}{2}}-\abs{x-z}^{-1+\frac{1}{2}}}\abs{\lapv b(y)} \abs{\lapv a(z)}}{\abs{x-y}^{1+s}} \ \d y\ \d z.
\]
For almost all $x,y,z$ we get from Lemma~\ref{lem:MultiplierEstimate} choosing $\delta:=\frac14+\frac s2$
\[
\frac{\abs{\abs{y-z}^{-1+\frac{1}{2}}-\abs{x-z}^{-1+\frac{1}{2}}}}{\abs{x-y}^{1+s}} \leq 
\brac{\abs{x-z}\abs{x-y}}^{-1+\fracm{4}-\frac{s}{2}}
+ \brac{\abs{y-z}\abs{x-y}}^{-1+\fracm{4}-\frac{s}{2}},
\]
which implies, using again~\eqref{eq:fraclap-}
\[
 \abs{II(x)} \aleq \lapms{\frac{1-2s}{4}} \abs{\lapv a}\ \lapms{\frac{1-2s}{4}} \abs{\lapv b} + \lapms{\frac{1-2s}{4}} \brac{\abs{\lapv b}\ \lapms{\frac{1-2s}{4}} \abs{\lapv a}}.
\]

By symmetry a respective estimate holds also for the term $III$.
Applying Lemma~\ref{lem:LowerOrderProducts.v2} one concludes.
\end{proof}



\begin{proof}[Proof of Lemma~\ref{lem:CriticalTerm}]
 Again the proof relies on quasi-locality. First we decompose
 \begin{align*}
  \left|\int \lapv  g_i' \omega_{ij} H_{1/2}(g'_j,\phi) \right|
  &\aleq  \left|\; \int_{B_{\Lambda^{1/2}r}}  \lapv  g_i' \omega_{ij} H_{1/2}(g'_j,\phi) \right|
  + \sum_{k=1}^\infty \left|\;
  \int_{A^k_{\Lambda^{1/2}, r}} 
  \lapv  g_i' \omega_{ij} H_{1/2}(g'_j,\phi) \right|\\
  &\aleq \|\lapv g'\|_{(2,\infty),B_{\Lambda^{1/2}r}} 
  \|H_{1/2}(g'_j,\phi)\|_{(2,1), B_{\Lambda^{1/2}r}} \\ 
  & \quad +
  \sum_{k=1}^\infty \|\lapv  g'\|_{(2,\infty),A_{\Lambda^{1/2},r}^{k}}
  \|H_{1/2}(g'_j,\phi)\|_{(2,1),A^k_{\Lambda^{1/2},r}}   .
 \end{align*}
  Using Lemma~\ref{lem:localizedHEst},
  the first summand can be estimated by %(for some uniform $\varepsilon > 0$)
 \begin{align*}
 \vrac{\lapms{s}\lapv \phi}_{(\frac{2}{1-2s},1)} \ \brac{
  \|\lapv g'\|_{(2,\infty),B_{\Lambda r}}^2 +  \|\lapv g'\|_{(2,\infty),B_{\Lambda r}}\ \Lambda^{-\theta} \sum_{l=1}^\infty 2^{-\theta l} \|\lapv  g'\|_{(2,\infty),B_{\Lambda 2^lr}}
  }.
 \end{align*}
 Applying the second part of Lemma~\ref{lem:localizedHEst}, the infinite sum can be estimated by
%  \begin{align*}
 \[ \Lambda^{-\theta}\| \lapms{s}  \lapv \phi\|_{(\frac 2 {1-2 s},1)}\ \|\lapv g'\|_{(2,\infty)} \ \sum_{k=2}^\infty 
  2^{-\theta k} \|\lapv  g'\|_{(2,\infty),B_{2^k\Lambda r}}.\]
% Since 
% $\| \lapms{\eps}  |\lapv \phi|\|_{(\frac 2 {1-2 \eps},1)}  
%  \aleq \|\lapv \phi\|_{(2,1)}$ by Sobolev's inequality, the 
%estimates above prove the claim.
%  \end{align*}
\end{proof}

% 
% \begin{lemma}\label{pr:localization}
% We have for $s \in (0,\frac{1}{2})$, for any $\Lambda > 2$,
% \[
%  \vrac{\lapms{s} f}_{(\frac{2}{1-2s},\infty),B_r} 
% \aleq \vrac{\chi_{B_{\Lambda r}}f}_{(2,\infty),B_r} + \sum_{k=1}^\infty 
% (\Lambda 2^k)^{s-\frac 1 2}\vrac{f}_{(2,\infty),A^k_{\Lambda,r}}.
% \]
% Here, we set
% \[
%  A^k_{\Lambda,r} := B_{2^{k} \Lambda r} \backslash B_{2^{k-1} \Lambda r}.
% \]
% \end{lemma}
% 
% \begin{proof}
% Using the triangle inequality we first get
% \begin{equation*}
%  \vrac{\lapms{s} f}_{(\frac{2}{1-2s},\infty),B_r} \aleq
% \vrac{\lapms{s} \chi_{B_{\Lambda r}}f}_{(\frac{2}{1-2s},\infty),B_r}
% + \sum_{k=1}^\infty \vrac{\lapms{s} \chi_{\eta^k_{\Lambda}}f}_{(\frac{2}{1-2s},\infty),B_r}
% \end{equation*}
% 
% By Sobolev's inequality, the first term can be estimated by
% \[
%  \vrac{\lapms{s} \chi_{B_{\Lambda r}}f}_{(\frac{2}{1-2s},\infty),B_r} \aleq \vrac{f}_{(2,\infty),B_{\Lambda r}}.
% \]
% For the second term we get using Lemma~\ref{lem:QuasiLocality} that
% \begin{align*}
%   \vrac{\lapms{s} \chi_{\etA^k_{\Lambda,r}}f}_{(\frac{2}{1-2s},\infty),B_r} &
%   \aleq (\Lambda 2^k r)^{-1+s} r^{1/2- s}(\Lambda r 2^k)^{1/2}\|f\|_{2,\infty; A^k_{\Lambda,r}}
%   \aleq (\Lambda 2^k)^{-1/2+s}
%  \end{align*}
% which finishes the proof.
% \end{proof}
% 
% \begin{remark}
%  Note that $\displaystyle\sum_{l=1}^\infty\Lambda^{-\theta l}\norm{u}_{(p,q),B_{\Lambda^kr}}
%  \approx\sum_{l=1}^\infty\tilde\Lambda^{-\tilde \theta l}\norm{u}_{(p,q),B_{\Lambda^kr}}$.
% \end{remark}



\subsection{Iteration lemma}

In order to prove Dirichlet growth, we need an iteration lemma
whose proof is based on the technique presented in~\cite[\href{http://www.math.ucla.edu/~tao/254a.1.01w/notes4.dvi}{notes4}, p.~11]{TaoLectures}.
The statement is also similar to the corresponding one appearing in~\cite{DR11Sphere}.
One should see this as a generalized version of De Giorgi's Iteration Lemma, cf., e.g., \cite{GiaquintaMI}.

\begin{lemma}[Iteration lemma] \label{lem:IterationLemma}
 Let $C<\infty$ and $\theta >0 $ be given.

 If $b_k\geq 0$, $k \in \mathbb N_0$, satisfy
 \begin{equation} \label{eq:Iteration}
  b_{k+m} \leq \varepsilon b_{k} +
    C\left(2^{-\theta (k+m)} + 2^{-\theta m}\sum_{l=1}^k 2^{-\theta l} b_{k-l}\right)  
 \end{equation}
 for all $k \in \mathbb N_0$ and $\varepsilon>0$ is small enough, $m$ is 
 big enough, then
 \begin{equation*}
  b_k \aleq 2^{-\tilde \theta k} 
 \end{equation*}
 for all $k\in \mathbb N_0$ with $\tilde \theta = \theta /2$.
\end{lemma}

\begin{proof}
 We will prove that
 \begin{equation} \label{eq:sumConverges}
   \sum_{k=0}^\infty 2^{\tilde \theta k} b_k \leq C\sum_{l=0}^m 2^{\tilde \theta l} b_l+ C.
 \end{equation}
 Especially, the infinite sum converges and hence the summands
 are a null series, which proves the lemma.

  Multiplying Equation~\eqref{eq:Iteration} with $2^{\tilde \theta k}$,
  and summing over $k$ we get 
 \begin{align*}
    2^{-\tilde \theta m} \sum_{k=0}^\infty {2}^{\tilde \theta (k+m)} b_{k+m}
     &\leq \varepsilon \sum_{k=0}^\infty 2^{\tilde \theta k} b_{k} 
     + C \sum_{k=0}^\infty 2^{-\theta (k+m)} 2^{\tilde \theta k} 
     + C 2^{-\theta m}\sum_{k=0}^\infty \sum_{l=1}^k 2^{-\theta l} 2^{\tilde \theta k} b_{k-l}\\
     &\leq \varepsilon \sum_{k=0}^\infty  2^{\tilde \theta k} b_{k}
     + C 2^{-\theta m} \sum_{k'=0}^\infty \sum_{l=1}^\infty 2^{-\theta l} 2^{\tilde \theta (k'+l)} b_{k'} 
     + C 2^{-\theta m} \\
     &\leq \varepsilon \sum_{k=0}^\infty 2^{\tilde \theta k} b_{k}
     + \tilde C 2^{-\theta m} \sum_{k'=0}^\infty 2^{\tilde \theta k'} b_{k'}
     + C 2^{-\theta m}.
 \end{align*}
  Assuming that $m$ is so large that
  $\tilde C 2^{-m(\theta  -\tilde \theta)} = \tilde C 2^{-\frac{m\theta}2} < 1/4$ and $0<\varepsilon<\frac 1 {4} 2^{-\theta m}$
  we get
 \begin{align*}
   \sum_{k=0}^\infty 2^{\tilde \theta k} b_k 
  \leq \frac 1 2 \sum_{k=0}^\infty 2^{\tilde \theta k} b_k  + \sum_{l=0}^m  2^{\tilde \theta l} b_l+ C
 \end{align*}
 and hence
 \begin{equation*}
  \sum_{k=0}^\infty 2^{\tilde \theta k} b_k \leq C \sum_{l=0}^m 2^{\tilde \theta l}b_l + C
 \end{equation*}
 if the infinite series converges. 
 If the sum is not known to converge, we apply the above argument to the cut-off series
 \begin{equation*}
  \begin{cases}
   \tilde b_k =b_k &\text{ if } k \le N, \\
   \tilde b_k = 0 &\text{ else}
  \end{cases}
 \end{equation*}
 and get the uniform bound
 $
  \sum_{k=0}^N  2^{\tilde \theta k} b_k \aleq \sum_{l=0}^m 2^{\tilde\theta l} b_l+ 1.
 $
  Letting $N \rightarrow \infty$ we get \eqref{eq:sumConverges}.
\end{proof}